%% u7-entries.tex -- Using InitechBase, Error Messages chapter (generated entries).
%% Source of the catalogue: os/samir/samir_main.c g_msg (l.128-279). 146 non-empty codes.

% 24 -- not raised: no source line sets ordinal 24 for the condition the message names (checked os/samir)
\begin{XbaseMessageEntry}[no command of this version raises this message]{24}{ALIAS name already in use.}
An alias chosen for a database is already the alias of another open database.
\end{XbaseMessageEntry}

% 13 -- os/samir/cmd/flow.c:1150 (SELECT <alias> with no such alias; a work-area number out of range is 17). Observed in the replay: SELECT FRED
\begin{XbaseMessageEntry}{13}{ALIAS not found.}
\cmd{SELECT} named an alias that is not the alias of any open database.

Use the alias given in \cmd{USE} \kw{ALIAS}, or the name of the database, or select the work area by its number.
\end{XbaseMessageEntry}

% 97 -- not raised: no source line sets ordinal 97 for the condition the message names (checked os/samir)
\begin{XbaseMessageEntry}[no command of this version raises this message]{97}{Bad read or illegal printer file.}
A read failed, or a printer file is not valid.
\end{XbaseMessageEntry}

% 38 -- not raised: no source line sets ordinal 38 for the condition the message names (checked os/samir)
\begin{XbaseMessageEntry}[no command of this version raises this message]{38}{Beginning of file encountered.}
The record pointer was moved backward past the first record.
\end{XbaseMessageEntry}

% 147 -- not raised: no source line sets ordinal 147 for the condition the message names (checked os/samir)
\begin{XbaseMessageEntry}[no command of this version raises this message]{147}{Cannot append in column order.}
Records cannot be appended in column order.
\end{XbaseMessageEntry}

% 89 -- not raised: no source line sets ordinal 89 for the condition the message names (checked os/samir)
\begin{XbaseMessageEntry}[no command of this version raises this message]{89}{Cannot erase a file which is open.}
An attempt was made to erase a file that is open.
\end{XbaseMessageEntry}

% 139 -- not raised: no source line sets ordinal 139 for the condition the message names (checked os/samir)
\begin{XbaseMessageEntry}[no command of this version raises this message]{139}{Cannot JOIN a file with itself.}
A database cannot be joined with itself.
\end{XbaseMessageEntry}

% 17 -- os/samir/samir_main.c:465 (USE: no selected area); os/samir/cmd/workarea.c:1334 (wa_set_filter on a closed area)
\begin{XbaseMessageEntry}{17}{Cannot select requested database.}
\cmd{SET FILTER TO} was given a condition while the selected work area holds no open database.

Open a database in the selected work area with \cmd{USE}, or select the area that holds one, and repeat the command.
\end{XbaseMessageEntry}

% 111 -- os/samir/cmd/mutate.c:124,618,833,848,895,971,1073 (-DBF_ERR_IO)
\begin{XbaseMessageEntry}{111}{Cannot write to a read-only file.}
A command that changes a database, \cmd{REPLACE}, \cmd{APPEND BLANK}, \cmd{DELETE}, \cmd{RECALL}, \cmd{PACK} or \cmd{ZAP}, found that the file is not open for writing.

Open the database with \cmd{USE} on a diskette that is not write-protected and repeat the command.
\end{XbaseMessageEntry}

% 57 -- os/samir/core/fn_builtins.c:316
\begin{XbaseMessageEntry}{57}{CHR() : Out of range.}
The argument of \fn{CHR()} was less than zero or greater than 255.

Give \fn{CHR()} a code from 0 through 255.
\end{XbaseMessageEntry}

% 66 -- not raised: no source line sets ordinal 66 for the condition the message names (checked os/samir)
\begin{XbaseMessageEntry}[no command of this version raises this message]{66}{CMDSET():}
A message introduces a command-set diagnostic.
\end{XbaseMessageEntry}

% 76 -- os/samir/core/eval.c:189 (XBEE_CONCAT_MINUS, result over 254 characters). Observed in the replay: X = SPACE(200), ? X - X
\begin{XbaseMessageEntry}{76}{- : Concatenated string too large.}
Two character strings joined with \texttt{-} give a string longer than 254 characters.

Shorten the strings, for example with \fn{TRIM()} or \fn{SUBSTR()}, before joining them.
\end{XbaseMessageEntry}

% 77 -- os/samir/core/eval.c:152 (XBEE_CONCAT_PLUS, result over 254 characters). Observed in the replay: X = SPACE(200), ? X + X
\begin{XbaseMessageEntry}{77}{+ : Concatenated string too large.}
Two character strings joined with \texttt{+} give a string longer than 254 characters.

Shorten the strings, for example with \fn{TRIM()} or \fn{SUBSTR()}, before joining them.
\end{XbaseMessageEntry}

% 42 -- os/samir/cmd/query.c:1218
\begin{XbaseMessageEntry}{42}{CONTINUE without LOCATE.}
\cmd{CONTINUE} was issued in a work area for which no \cmd{LOCATE} has been given.

Issue \cmd{LOCATE FOR} with a condition in the same work area, then \cmd{CONTINUE}.
\end{XbaseMessageEntry}

% 72 -- not raised: no source line sets ordinal 72 for the condition the message names (checked os/samir)
\begin{XbaseMessageEntry}[no command of this version raises this message]{72}{could not be opened.}
A file named in a message could not be opened.
\end{XbaseMessageEntry}

% 44 -- not raised: no source line sets ordinal 44 for the condition the message names (checked os/samir)
\begin{XbaseMessageEntry}[no command of this version raises this message]{44}{Cyclic relation.}
A chain of relations between databases leads back to the database it began with.
\end{XbaseMessageEntry}

% 122 -- not raised: no source line sets ordinal 122 for the condition the message names (checked os/samir)
\begin{XbaseMessageEntry}[no command of this version raises this message]{122}{Data Catalog has not been established.}
No data catalog has been set up.
\end{XbaseMessageEntry}

% 9 -- os/samir/core/eval.c:345,428,475,503,511,526,534,562,595,602 (operators); os/samir/core/fn_builtins.c:191 and following (function arguments); os/samir/cmd/mutate.c:605,617 (REPLACE); os/samir/cmd/proc.c:308 (ACCEPT/INPUT prompt); os/samir/cmd/report.c:251,278,474
\begin{XbaseMessageEntry}{9}{Data type mismatch.}
An expression combined operands of different types, for example a character value and a numeric value with the plus sign, or gave a function an argument of the wrong type. \cmd{REPLACE} reports the message when the value does not match the type of the field, or the field does not exist. \cmd{ACCEPT} and \cmd{INPUT} report it when the prompt is not a character expression, and \cmd{REPORT FORM} reports it when a heading or group expression has the wrong type.

Convert one operand with \fn{STR()}, \fn{VAL()}, \fn{DTOC()} or \fn{CTOD()}, or correct the field name or the value in the command.
\end{XbaseMessageEntry}

% 71 -- not raised: no source line sets ordinal 71 for the condition the message names (checked os/samir)
\begin{XbaseMessageEntry}[no command of this version raises this message]{71}{, Database in Use:}
The text introduces the name of the database in use in a status display.
\end{XbaseMessageEntry}

% 131 -- not raised: no source line sets ordinal 131 for the condition the message names (checked os/samir)
\begin{XbaseMessageEntry}[no command of this version raises this message]{131}{Database is encrypted.}
A database is encrypted.
\end{XbaseMessageEntry}

% 26 -- os/samir/cmd/query.c:1302,1307 (SEEK/FIND without a master index); query.c:931
\begin{XbaseMessageEntry}{26}{Database is not indexed.}
\cmd{SEEK} or \cmd{FIND} was issued for a database that has no master index in use.

Open the database with an index, \cmd{USE STAFF INDEX STAFFNAM}, or select an index with \cmd{SET ORDER TO}, and repeat the search.
\end{XbaseMessageEntry}

% 41 -- not raised: no source line sets ordinal 41 for the condition the message names (checked os/samir)
\begin{XbaseMessageEntry}[no command of this version raises this message]{41}{.DBT file cannot be opened.}
The memo file that belongs to a database cannot be opened.
\end{XbaseMessageEntry}

% 56 -- os/samir/cmd/mutate.c:130 (-DBF_ERR_NOSPC); os/samir/samir_main.c:426
\begin{XbaseMessageEntry}{56}{Disk full when writing file:}
A command that changes a database file, such as \cmd{APPEND BLANK}, \cmd{REPLACE}, \cmd{DELETE} or \cmd{PACK}, could not write all of its data because the diskette is full. The name of the file follows the message.

Free space on the diskette by erasing files that are no longer needed, or use another diskette, and repeat the command.
\end{XbaseMessageEntry}

% 103 -- os/samir/cmd/flow.c:1460 fl_catalog (m[8]). Observed in the replay: a procedure that calls itself with DO
\begin{XbaseMessageEntry}{103}{DOs nested too deep.}
Programs and procedures call one another with \cmd{DO} more deeply than InitechBase allows; usually a procedure calls itself, directly or through another, without end.

End each procedure with \cmd{RETURN} to its caller instead of calling the program that called it.
\end{XbaseMessageEntry}

% 4 -- not raised: no source line sets ordinal 4 for the condition the message names (checked os/samir)
\begin{XbaseMessageEntry}[no command of this version raises this message]{4}{End of file encountered.}
The record pointer was moved, or a read was made, beyond the end of the file.
\end{XbaseMessageEntry}

% 51 -- not raised: no source line sets ordinal 51 for the condition the message names (checked os/samir)
\begin{XbaseMessageEntry}[no command of this version raises this message]{51}{End of file or error on keyboard input.}
The end of the keyboard input was reached, or the keyboard returned an error.
\end{XbaseMessageEntry}

% 145 -- not raised: no source line sets ordinal 145 for the condition the message names (checked os/samir)
\begin{XbaseMessageEntry}[no command of this version raises this message]{145}{Error in configuration value.}
A configuration value is in error.
\end{XbaseMessageEntry}

% 85 -- not raised: no source line sets ordinal 85 for the condition the message names (checked os/samir)
\begin{XbaseMessageEntry}[no command of this version raises this message]{85}{Error:}
The text introduces the description of an error.
\end{XbaseMessageEntry}

% 67 -- os/samir/cmd/workarea.c:1401,1402 (g_lex_msg / g_parse_msg: the evaluator's token or node pool is exhausted). Observed in the replay: ? with 100 nested parentheses, and ? 1+1+... with 120 terms
\begin{XbaseMessageEntry}{67}{EVAL work area overflow}
An expression is too long or too deeply nested to be evaluated.

Break the calculation into steps, storing each partial result in a memory variable.
\end{XbaseMessageEntry}

% 110 -- not raised: no source line sets ordinal 110 for the condition the message names (checked os/samir)
\begin{XbaseMessageEntry}[no command of this version raises this message]{110}{Exclusive open of file is required.}
A file must be opened for exclusive use before the operation can be done.
\end{XbaseMessageEntry}

% 80 -- not raised: no source line sets ordinal 80 for the condition the message names (checked os/samir)
\begin{XbaseMessageEntry}[no command of this version raises this message]{80}{***Execution error on}
The text introduces the name of the command in which an execution error occurred.
\end{XbaseMessageEntry}

% 73 -- not raised: no source line sets ordinal 73 for the condition the message names (checked os/samir)
\begin{XbaseMessageEntry}[no command of this version raises this message]{73}{^^ Expected ON or OFF.}
A command that takes ON or OFF was given another word.
\end{XbaseMessageEntry}

% 48 -- os/samir/cmd/mutate.c d_sort, d_total (field not in the table); lane samirdml 14a8fc75
\begin{XbaseMessageEntry}{48}{Field not found.}
\cmd{SORT} or \cmd{TOTAL} named a field that is not a field of the database in use.  Check the spelling of the field name.
\end{XbaseMessageEntry}

% 141 -- not raised: no source line sets ordinal 141 for the condition the message names (checked os/samir)
\begin{XbaseMessageEntry}[no command of this version raises this message]{141}{Fields list too complicated.}
A list of fields is too complicated to process.
\end{XbaseMessageEntry}

% 7 -- os/samir/cmd/mutate.c d_destination (SAFETY ON, target exists: INDEX, SORT, TOTAL, RENAME, COPY FILE); lane samirdml 14a8fc75
\begin{XbaseMessageEntry}{7}{File already exists.}
\cmd{INDEX}, \cmd{SORT}, \cmd{TOTAL} or \cmd{RENAME} was to create a file whose name is already on the diskette, and SET SAFETY is ON.  Choose another name, \cmd{ERASE} the old file, or SET SAFETY OFF.
\end{XbaseMessageEntry}

% 1 -- os/samir/samir_main.c:539 (USE: -DBF_ERR_NOENT); samir_main.c:673 (DO program read fault); os/samir/cmd/report.c:322 (REPORT FORM: form file absent)
\begin{XbaseMessageEntry}{1}{File does not exist.}
The \cmd{USE} command could not find the database named, or an index file named in its \kw{INDEX} clause, on the diskette. \cmd{REPORT FORM} reports the same message when the report form file named does not exist. A program file that cannot be read in full while \cmd{DO} loads it is reported the same way.

Check the spelling of the file name and the diskette that is in use. When no extension is typed, \cmd{USE} supplies \file{.DBF} and an index name receives \file{.NDX}, and \cmd{REPORT FORM} supplies \file{.FRM}.
\end{XbaseMessageEntry}

% 49 -- not raised: no source line sets ordinal 49 for the condition the message names (checked os/samir)
\begin{XbaseMessageEntry}[no command of this version raises this message]{49}{File has been deleted.}
A file has been deleted and cannot be used.
\end{XbaseMessageEntry}

% 3 -- os/samir/samir_main.c:521 (USE: wa_name_is_open); samir_main.c:539 (-WA_ERR_ALREADY_OPEN)
\begin{XbaseMessageEntry}{3}{File is already open.}
The \cmd{USE} command named a database file that is already open in another work area.

Select the work area that holds the file with \cmd{SELECT}, or close it with \cmd{CLOSE} before opening it again.
\end{XbaseMessageEntry}

% 108 -- not raised: no source line sets ordinal 108 for the condition the message names (checked os/samir)
\begin{XbaseMessageEntry}[no command of this version raises this message]{108}{File is in use by another.}
A file is in use by another user.
\end{XbaseMessageEntry}

% 29 -- os/samir/samir_main.c:540 (USE access fault); os/samir/cmd/mutate.c:131 (write refused); os/samir/cmd/report.c:110,322,482; os/samir/cmd/query.c:822,1333; os/samir/cmd/workarea.c:1363
\begin{XbaseMessageEntry}{29}{File is not accessible.}
A file could not be opened, read or created. \cmd{REPORT FORM} reports the message when the \kw{TO FILE} file cannot be created or written, and a command that changes a database reports it when the diskette refuses access to the file.

Check that the diskette is in the drive and not full of damaged sectors, and that the path in the file name exists.
\end{XbaseMessageEntry}

% 91 -- not raised: no source line sets ordinal 91 for the condition the message names (checked os/samir)
\begin{XbaseMessageEntry}[no command of this version raises this message]{91}{File was not LOADed.}
A module was used that had not been loaded.
\end{XbaseMessageEntry}

% 150 -- not raised: no source line sets ordinal 150 for the condition the message names (checked os/samir)
\begin{XbaseMessageEntry}[no command of this version raises this message]{150}{Help text not found.}
No help text was found for the topic requested.
\end{XbaseMessageEntry}

% 68 -- not raised: no source line sets ordinal 68 for the condition the message names (checked os/samir)
\begin{XbaseMessageEntry}[no command of this version raises this message]{68}{Illegal opcode}
An operation code that is not valid was met.
\end{XbaseMessageEntry}

% 46 -- not raised: no source line sets ordinal 46 for the condition the message names (checked os/samir)
\begin{XbaseMessageEntry}[no command of this version raises this message]{46}{Illegal value.}
A value was given that the command does not accept.
\end{XbaseMessageEntry}

% 114 -- os/samir/samir_main.c:555
\begin{XbaseMessageEntry}{114}{Index damaged.  REINDEX should be done before using data.}
\cmd{USE} opened the database and its indexes but could not position the file at the first record in index order, which indicates that an index is damaged.

Rebuild the index from the database and open the files again.
\end{XbaseMessageEntry}

% 112 -- not raised: no source line sets ordinal 112 for the condition the message names (checked os/samir)
\begin{XbaseMessageEntry}[no command of this version raises this message]{112}{Index expression is too big (220 char maximum).}
An index expression is longer than 220 characters.
\end{XbaseMessageEntry}

% 19 -- not raised: no source line sets ordinal 19 for the condition the message names (checked os/samir)
\begin{XbaseMessageEntry}[no command of this version raises this message]{19}{Index file does not match database.}
An index file was opened with a database other than the one it was built for.
\end{XbaseMessageEntry}

% 113 -- not raised: no source line sets ordinal 113 for the condition the message names (checked os/samir)
\begin{XbaseMessageEntry}[no command of this version raises this message]{113}{Index interrupted.  Index will be deleted if not completed.}
Building an index was interrupted, and the incomplete index will be deleted.
\end{XbaseMessageEntry}

% 23 -- os/samir/cmd/mutate.c d_sort (character key over 100); lane samirdml 14a8fc75
\begin{XbaseMessageEntry}{23}{Index is too big (100 char maximum).}
\cmd{SORT} was asked to order the records on a character field longer than 100 characters.
\end{XbaseMessageEntry}

% 43 -- os/samir/samir_main.c:657 (DO buffer); os/samir/cmd/report.c:236,273 (cell text pool)
\begin{XbaseMessageEntry}{43}{Insufficient memory.}
InitechBase had no room left to hold a program file that \cmd{DO} was loading, or the text of one record of a report.

Close databases that are not needed, shorten the program, or narrow the report columns, and repeat the command.
\end{XbaseMessageEntry}

% 64 -- not raised: no source line sets ordinal 64 for the condition the message names (checked os/samir)
\begin{XbaseMessageEntry}[no command of this version raises this message]{64}{Internal error:}
An internal inconsistency was detected.
\end{XbaseMessageEntry}

% 81 -- not raised: no source line sets ordinal 81 for the condition the message names (checked os/samir)
\begin{XbaseMessageEntry}[no command of this version raises this message]{81}{Invalid date.}
A date is not a valid date.
\end{XbaseMessageEntry}

% 118 -- not raised: no source line sets ordinal 118 for the condition the message names (checked os/samir)
\begin{XbaseMessageEntry}[no command of this version raises this message]{118}{Invalid DIF character.}
A character in a DIF file is not valid.
\end{XbaseMessageEntry}

% 115 -- not raised: no source line sets ordinal 115 for the condition the message names (checked os/samir)
\begin{XbaseMessageEntry}[no command of this version raises this message]{115}{Invalid DIF File Header.}
The header of an exchange file in DIF form is not valid.
\end{XbaseMessageEntry}

% 117 -- not raised: no source line sets ordinal 117 for the condition the message names (checked os/samir)
\begin{XbaseMessageEntry}[no command of this version raises this message]{117}{Invalid DIF type indicator.}
A type indicator in a DIF file is not valid.
\end{XbaseMessageEntry}

% 116 -- not raised: no source line sets ordinal 116 for the condition the message names (checked os/samir)
\begin{XbaseMessageEntry}[no command of this version raises this message]{116}{Invalid DIF vector - DBF field mismatch.}
A vector of a DIF file does not match the fields of the database.
\end{XbaseMessageEntry}

% 99 -- not raised: no source line sets ordinal 99 for the condition the message names (checked os/samir)
\begin{XbaseMessageEntry}[no command of this version raises this message]{99}{Invalid DOS SET option.}
An operating system option given to the SET command is not valid.
\end{XbaseMessageEntry}

% 11 -- os/samir/core/fn_builtins.c:190,237,280,293,313 and following (argument count); os/samir/core/eval.c:673-703 (IIF)
\begin{XbaseMessageEntry}{11}{Invalid function argument.}
A function was called with too few or too many arguments, or the argument list of \fn{IIF()} was not three expressions.

Give the function the number of arguments its entry in this volume shows.
\end{XbaseMessageEntry}

% 31 -- os/samir/core/eval.c:664; os/samir/core/fn_builtins.c:2094,2176
\begin{XbaseMessageEntry}{31}{Invalid function name.}
An expression named a function that InitechBase does not provide.

Check the spelling of the function name against the functions listed in this volume.
\end{XbaseMessageEntry}

% 106 -- os/samir/cmd/set.c:372 (SET ORDER TO <n> with no n-th open index). Observed in the replay: USE STAFF INDEX STAFFNAM, SET ORDER TO 9
\begin{XbaseMessageEntry}{106}{Invalid index number.}
\cmd{SET ORDER TO} gave the number of an index file that is not open: fewer index files are open in the work area than the number.

Give a number from 1 to the number of open index files, or 0 for the order of the records.
\end{XbaseMessageEntry}

% 107 -- not raised: no source line sets ordinal 107 for the condition the message names (checked os/samir)
\begin{XbaseMessageEntry}[no command of this version raises this message]{107}{Invalid operator.}
An operator is not valid for the operands it was given.
\end{XbaseMessageEntry}

% 123 -- not raised: no source line sets ordinal 123 for the condition the message names (checked os/samir)
\begin{XbaseMessageEntry}[no command of this version raises this message]{123}{Invalid printer port.}
A printer port is not valid.
\end{XbaseMessageEntry}

% 124 -- not raised: no source line sets ordinal 124 for the condition the message names (checked os/samir)
\begin{XbaseMessageEntry}[no command of this version raises this message]{124}{Invalid printer redirection.}
A printer redirection is not valid.
\end{XbaseMessageEntry}

% 120 -- not raised: no source line sets ordinal 120 for the condition the message names (checked os/samir)
\begin{XbaseMessageEntry}[no command of this version raises this message]{120}{Invalid SYLK file dimension bounds.}
The dimension bounds of a SYLK file are not valid.
\end{XbaseMessageEntry}

% 121 -- not raised: no source line sets ordinal 121 for the condition the message names (checked os/samir)
\begin{XbaseMessageEntry}[no command of this version raises this message]{121}{Invalid SYLK file format.}
A SYLK file is not in a valid format.
\end{XbaseMessageEntry}

% 119 -- not raised: no source line sets ordinal 119 for the condition the message names (checked os/samir)
\begin{XbaseMessageEntry}[no command of this version raises this message]{119}{Invalid SYLK file header.}
The header of an exchange file in SYLK form is not valid.
\end{XbaseMessageEntry}

% 83 -- not raised: no source line sets ordinal 83 for the condition the message names (checked os/samir)
\begin{XbaseMessageEntry}[no command of this version raises this message]{83}{is not a dBASE command.}
The text follows a word that is not the name of a command.
\end{XbaseMessageEntry}

% 86 -- not raised: no source line sets ordinal 86 for the condition the message names (checked os/samir)
\begin{XbaseMessageEntry}[no command of this version raises this message]{86}{^--- Keyword not found.}
A keyword that a command requires was not found.
\end{XbaseMessageEntry}

% 54 -- os/samir/cmd/report.c lbl_run (not a 1034-byte form, or sizes out of range); lane samirdml 14a8fc75
\begin{XbaseMessageEntry}{54}{Label file invalid.}
The file named in \cmd{LABEL FORM} is not a dBASE III PLUS label form, or its label sizes are out of range.
\end{XbaseMessageEntry}

% 18 -- os/samir/samir_main.c:650 (DO: program file of 8,192 bytes or more); os/samir/cmd/report.c:430 (command text of 256 or more characters)
\begin{XbaseMessageEntry}{18}{Line exceeds maximum of 254 characters.}
\cmd{DO} was given a program file of 8,192 bytes or more, which is too large to load. \cmd{REPORT FORM} reports the same message for a command whose text is too long.

Divide the program into smaller files that call one another with \cmd{DO}, and keep each command line within 254 characters.
\end{XbaseMessageEntry}

% 100 -- not raised: no source line sets ordinal 100 for the condition the message names (checked os/samir)
\begin{XbaseMessageEntry}[no command of this version raises this message]{100}{Lock failed, but not because of previous lock.}
A lock request failed for a reason other than an existing lock.
\end{XbaseMessageEntry}

% 58 -- not raised: no source line sets ordinal 58 for the condition the message names (checked os/samir)
\begin{XbaseMessageEntry}[no command of this version raises this message]{58}{LOG() : Zero or negative.}
The argument given to a logarithm function is zero or negative.
\end{XbaseMessageEntry}

% 149 -- not raised: no source line sets ordinal 149 for the condition the message names (checked os/samir)
\begin{XbaseMessageEntry}[no command of this version raises this message]{149}{Master catalog is empty.}
The master catalog is empty.
\end{XbaseMessageEntry}

% 146 -- not raised: no source line sets ordinal 146 for the condition the message names (checked os/samir)
\begin{XbaseMessageEntry}[no command of this version raises this message]{146}{Maximum path length exceeded.}
A path is longer than the maximum allowed.
\end{XbaseMessageEntry}

% 137 -- not raised: no source line sets ordinal 137 for the condition the message names (checked os/samir)
\begin{XbaseMessageEntry}[no command of this version raises this message]{137}{Maximum record length exceeded.}
The record length is greater than the maximum allowed.
\end{XbaseMessageEntry}

% 55 -- os/samir/cmd/flow.c fl_memfile (RESTORE framing checked first); lane samirdml 14a8fc75
\begin{XbaseMessageEntry}{55}{Memory Variable file is invalid.}
The file named in \cmd{RESTORE FROM} is not a memory file written by \cmd{SAVE}.  No memory variable was changed.
\end{XbaseMessageEntry}

% 96 -- os/samir/cmd/flow.c:1460 fl_catalog (m[5], -INTERP_ERR_STRUCT in a program) from flow.c:1544 (a lone ENDDO). Observed in the replay: a program whose first line is ENDDO
\begin{XbaseMessageEntry}{96}{Mismatched DO WHILE and ENDDO.}
A program has an \cmd{ENDDO} without a \cmd{DO WHILE} before it, or a \cmd{DO WHILE} without its \cmd{ENDDO}.

Give every \cmd{DO WHILE} of the program its own \cmd{ENDDO}, and check that the \cmd{IF} and \cmd{DO CASE} groups inside it are complete.
\end{XbaseMessageEntry}

% 87 -- not raised: no source line sets ordinal 87 for the condition the message names (checked os/samir)
\begin{XbaseMessageEntry}[no command of this version raises this message]{87}{NDX(): Invalid index number.}
The index number given to an index function does not name an open index.
\end{XbaseMessageEntry}

% 148 -- not raised: no source line sets ordinal 148 for the condition the message names (checked os/samir)
\begin{XbaseMessageEntry}[no command of this version raises this message]{148}{Network server busy.}
The network server is busy.
\end{XbaseMessageEntry}

% 52 -- os/samir/cmd/report.c:466; os/samir/core/fn_builtins.c:1868 (need_dbcur)
\begin{XbaseMessageEntry}{52}{No database is in USE.}
\cmd{REPORT FORM} was issued while the selected work area holds no open database.

Open the database with \cmd{USE} and repeat the command.
\end{XbaseMessageEntry}

% 47 -- not raised: no source line sets ordinal 47 for the condition the message names (checked os/samir)
\begin{XbaseMessageEntry}[no command of this version raises this message]{47}{No fields to process.}
A command that works on fields found no fields to work on.
\end{XbaseMessageEntry}

% 138 -- not raised: no source line sets ordinal 138 for the condition the message names (checked os/samir)
\begin{XbaseMessageEntry}[no command of this version raises this message]{138}{No fields were found to copy.}
No fields were found to copy.
\end{XbaseMessageEntry}

% 14 -- not raised: no source line sets ordinal 14 for the condition the message names (checked os/samir)
\begin{XbaseMessageEntry}[no command of this version raises this message]{14}{No find.}
A search for a key found no matching record.
\end{XbaseMessageEntry}

% 93 -- os/samir/cmd/proc.c:676
\begin{XbaseMessageEntry}{93}{No PARAMETER statement found.}
\cmd{DO} with a \kw{WITH} list called a procedure that contains no \cmd{PARAMETERS} statement.

Begin the procedure with a \cmd{PARAMETERS} statement that names a variable for each value passed, or remove the \kw{WITH} list.
\end{XbaseMessageEntry}

% 45 -- not raised: no source line sets ordinal 45 for the condition the message names (checked os/samir)
\begin{XbaseMessageEntry}[no command of this version raises this message]{45}{Not a Character expression.}
A character expression was required and the expression given has another type.
\end{XbaseMessageEntry}

% 15 -- os/samir/samir_main.c:540 (USE: any other open fault)
\begin{XbaseMessageEntry}{15}{Not a dBASE database.}
\cmd{USE} opened a file whose header is not that of a database file, or an index file that does not belong with the database.

Name a file that was created as a database. Do not give \cmd{USE} a report form, a program or a text file.
\end{XbaseMessageEntry}

% 37 -- os/samir/core/eval.c:566-609,685; os/samir/cmd/flow.c:476; os/samir/cmd/query.c:234; os/samir/cmd/mutate.c:250; os/samir/cmd/workarea.c:1339,1373; os/samir/cmd/report.c:369
\begin{XbaseMessageEntry}{37}{Not a Logical expression.}
A condition that must be logical produced a value of another type. This applies to the \kw{FOR} and \kw{WHILE} clauses, to \cmd{IF}, \cmd{DO WHILE}, \cmd{DO CASE} and \cmd{SET FILTER TO}, and to the operands of \kw{.AND.}, \kw{.OR.} and \kw{.NOT.}. A number or character value is never treated as true or false.

Write a comparison, for example \type{SALARY > 0}, or use a logical field or the logical constants \kw{.T.} and \kw{.F.}.
\end{XbaseMessageEntry}

% 27 -- not raised: no source line sets ordinal 27 for the condition the message names (checked os/samir)
\begin{XbaseMessageEntry}[no command of this version raises this message]{27}{Not a numeric expression.}
A numeric expression was required and the expression given has another type.
\end{XbaseMessageEntry}

% 98 -- not raised: no source line sets ordinal 98 for the condition the message names (checked os/samir)
\begin{XbaseMessageEntry}[no command of this version raises this message]{98}{Not a RunTime file.}
A file is not a run-time file.
\end{XbaseMessageEntry}

% 140 -- not raised: no source line sets ordinal 140 for the condition the message names (checked os/samir)
\begin{XbaseMessageEntry}[no command of this version raises this message]{140}{Not a valid PFS file.}
A file is not a valid PFS file.
\end{XbaseMessageEntry}

% 134 -- not raised: no source line sets ordinal 134 for the condition the message names (checked os/samir)
\begin{XbaseMessageEntry}[no command of this version raises this message]{134}{Not a valid QUERY file.}
A file is not a valid query file.
\end{XbaseMessageEntry}

% 127 -- not raised: no source line sets ordinal 127 for the condition the message names (checked os/samir)
\begin{XbaseMessageEntry}[no command of this version raises this message]{127}{Not a valid VIEW file.}
A file is not a valid view file.
\end{XbaseMessageEntry}

% 82 -- not raised: no source line sets ordinal 82 for the condition the message names (checked os/samir)
\begin{XbaseMessageEntry}[no command of this version raises this message]{82}{** Not Found **}
A search ended without finding the key sought.
\end{XbaseMessageEntry}

% 101 -- not raised: no source line sets ordinal 101 for the condition the message names (checked os/samir)
\begin{XbaseMessageEntry}[no command of this version raises this message]{101}{Not suspended.}
A program was resumed that had not been suspended.
\end{XbaseMessageEntry}

% 39 -- os/samir/core/eval.c:521 (division by zero)
\begin{XbaseMessageEntry}{39}{Numeric overflow (data was lost).}
An expression divided by zero. InitechBase reports the quotient as a numeric overflow.

Test the divisor before dividing, for example with \fn{IIF()}.
\end{XbaseMessageEntry}

% 90 -- not raised: no source line sets ordinal 90 for the condition the message names (checked os/samir)
\begin{XbaseMessageEntry}[no command of this version raises this message]{90}{Operation with Logical field invalid.}
An operation was attempted on a logical field that logical fields do not allow.
\end{XbaseMessageEntry}

% 34 -- os/samir/cmd/report.c:234 (rp_text: XB_M value)
\begin{XbaseMessageEntry}{34}{Operation with Memo field invalid.}
\cmd{REPORT FORM} met a memo field where a report column or heading needs a character, numeric or date value.

Edit the report form so that its columns use character, numeric, date or logical expressions only.
\end{XbaseMessageEntry}

% 78 -- not raised: no source line sets ordinal 78 for the condition the message names (checked os/samir)
\begin{XbaseMessageEntry}[no command of this version raises this message]{78}{^ or ** : Negative base, fractional exponent.}
A negative number was raised to a fractional power.
\end{XbaseMessageEntry}

% 21 -- os/samir/cmd/flow.c:344,413,430,440 (-INTERP_ERR_NOMEM) -> flow.c:469 fl_memvar_fail; flow.c:1052 (RESTORE). The OS pool is the III+ 6000 bytes (os/samir/pal/pal_milton.c, Using III+ U5-180). Observed in the replay: a program storing SPACE(250) in 250 memory variables
\begin{XbaseMessageEntry}{21}{Out of memory variable memory.}
A value could not be stored in a memory variable because the space for the values of memory variables is full.  It is also given by \cmd{RESTORE} when the variables of the memory file do not fit.

Release the memory variables that are no longer needed with \cmd{RELEASE}, or store shorter character values.
\end{XbaseMessageEntry}

% 22 -- os/samir/cmd/flow.c:408 (FLOW_MAX_MEMVARS 256, -INTERP_ERR_DEPTH) -> flow.c:469 fl_memvar_fail. Observed in the replay: a program creating 400 memory variables with &
\begin{XbaseMessageEntry}{22}{Out of memory variable slots.}
A new memory variable could not be created because 256 memory variables already exist.

Release the memory variables that are no longer needed with \cmd{RELEASE}, or use fewer of them.
\end{XbaseMessageEntry}

% 75 -- not raised: no source line sets ordinal 75 for the condition the message names (checked os/samir)
\begin{XbaseMessageEntry}[no command of this version raises this message]{75}{^--- Out of range.}
A value lies outside its permitted range.
\end{XbaseMessageEntry}

% 30 -- not raised: no source line sets ordinal 30 for the condition the message names (checked os/samir)
\begin{XbaseMessageEntry}[no command of this version raises this message]{30}{Position is off the screen.}
A screen row or column lies outside the screen.
\end{XbaseMessageEntry}

% 126 -- os/samir/cmd/report.c:461
\begin{XbaseMessageEntry}{126}{Printer is either not connected or turned off.}
\cmd{REPORT FORM} was given \kw{TO PRINT}. This version has no printer to send the report to.

Send the report to the screen, or to a text file with \kw{TO FILE}.
\end{XbaseMessageEntry}

% 125 -- not raised: no source line sets ordinal 125 for the condition the message names (checked os/samir)
\begin{XbaseMessageEntry}[no command of this version raises this message]{125}{Printer not ready.}
The printer is not ready.
\end{XbaseMessageEntry}

% 143 -- not raised: no source line sets ordinal 143 for the condition the message names (checked os/samir)
\begin{XbaseMessageEntry}[no command of this version raises this message]{143}{Query not valid for this environment.}
A query cannot be used in the present environment.
\end{XbaseMessageEntry}

% 109 -- not raised: no source line sets ordinal 109 for the condition the message names (checked os/samir)
\begin{XbaseMessageEntry}[no command of this version raises this message]{109}{Record is in use by another.}
A record is in use by another user.
\end{XbaseMessageEntry}

% 20 -- not raised: no source line sets ordinal 20 for the condition the message names (checked os/samir)
\begin{XbaseMessageEntry}[no command of this version raises this message]{20}{Record is not in index.}
A record of the database has no entry in the index in use.
\end{XbaseMessageEntry}

% 25 -- not raised: no source line sets ordinal 25 for the condition the message names (checked os/samir)
\begin{XbaseMessageEntry}[no command of this version raises this message]{25}{Record is not inserted.}
A record has been created but has not yet been placed in the database.
\end{XbaseMessageEntry}

% 130 -- not raised: no source line sets ordinal 130 for the condition the message names (checked os/samir)
\begin{XbaseMessageEntry}[no command of this version raises this message]{130}{Record is not locked.}
An operation requires a record to be locked and it is not.
\end{XbaseMessageEntry}

% 5 -- os/samir/cmd/query.c:1423 (GO/GOTO: wa_nav_goto fails); os/samir/cmd/report.c:387 (REPORT FORM RECORD clause)
\begin{XbaseMessageEntry}{5}{Record is out of range.}
\cmd{GO} or \cmd{GOTO} was given a record number that is less than one or greater than the number of records in the database. \cmd{REPORT FORM} reports the message when its \kw{RECORD} clause names a record that does not exist.

Test the record count with \fn{RECCOUNT()} and name a record number within it, or use \kw{TOP} or \kw{BOTTOM}.
\end{XbaseMessageEntry}

% 142 -- not raised: no source line sets ordinal 142 for the condition the message names (checked os/samir)
\begin{XbaseMessageEntry}[no command of this version raises this message]{142}{Relation record is in use by others.}
A related record is in use by other users.
\end{XbaseMessageEntry}

% 88 -- not raised: no source line sets ordinal 88 for the condition the message names (checked os/samir)
\begin{XbaseMessageEntry}[no command of this version raises this message]{88}{REPLICATE(): String too large.}
The result of \fn{REPLICATE()} would be longer than a character string may be.
\end{XbaseMessageEntry}

% 50 -- os/samir/cmd/report.c:327-355 (rp_load)
\begin{XbaseMessageEntry}{50}{Report file invalid.}
\cmd{REPORT FORM} read a file that is not a valid report form, or a report form that uses an option this version does not print.

Name a file that was made as a report form, and do not edit it with a program that does not preserve its layout.
\end{XbaseMessageEntry}

% 60 -- os/samir/core/fn_builtins.c:296
\begin{XbaseMessageEntry}{60}{SPACE() : Negative.}
The argument of \fn{SPACE()} was less than zero.

Give \fn{SPACE()} zero or a positive number.
\end{XbaseMessageEntry}

% 59 -- os/samir/core/fn_builtins.c:297,1203,1301,1565,1768
\begin{XbaseMessageEntry}{59}{SPACE() : Too large.}
\fn{SPACE()}, \fn{REPLICATE()}, \fn{STUFF()} or another function would have produced a character string longer than 254 characters.

Request a shorter result.
\end{XbaseMessageEntry}

% 61 -- not raised: no source line sets ordinal 61 for the condition the message names (checked os/samir)
\begin{XbaseMessageEntry}[no command of this version raises this message]{61}{SQRT() : Negative.}
The argument given to a square root function is negative.
\end{XbaseMessageEntry}

% 79 -- not raised: no source line sets ordinal 79 for the condition the message names (checked os/samir)
\begin{XbaseMessageEntry}[no command of this version raises this message]{79}{STORE : String too large.}
A character string given to \cmd{STORE} is longer than a memory variable may hold.
\end{XbaseMessageEntry}

% 63 -- os/samir/core/fn_builtins.c:383-385
\begin{XbaseMessageEntry}{63}{STR() : Out of range.}
The length or the number of decimal places given to \fn{STR()} was not usable: a length of less than one or more than 254, more than 15 decimal places, or too little room for the decimals.

Give a length that can hold the digits, the sign and the decimal point.
\end{XbaseMessageEntry}

% 33 -- not raised: no source line sets ordinal 33 for the condition the message names (checked os/samir)
\begin{XbaseMessageEntry}[no command of this version raises this message]{33}{Structure invalid.}
The structure of a database is not valid.
\end{XbaseMessageEntry}

% 102 -- not raised: no source line sets ordinal 102 for the condition the message names (checked os/samir)
\begin{XbaseMessageEntry}[no command of this version raises this message]{102}{STUFF():  String too large.}
The result of \fn{STUFF()} would be longer than a character string may be.
\end{XbaseMessageEntry}

% 62 -- os/samir/core/fn_builtins.c:247
\begin{XbaseMessageEntry}{62}{SUBSTR() : Start point out of range.}
The starting position given to \fn{SUBSTR()} was less than one.

Give a starting position of one or more.
\end{XbaseMessageEntry}

% 10 -- os/samir/samir_main.c:481,506,515 (USE); os/samir/cmd/set.c:600,635 (SET DELETED/HEADING/CONSOLE/STATUS); os/samir/cmd/query.c:789,1323,1486; os/samir/cmd/mutate.c:1099,1111; os/samir/cmd/workarea.c:1335; os/samir/cmd/report.c:309,384
\begin{XbaseMessageEntry}{10}{Syntax error.}
A command was typed with a required word missing or an operand that is too long. Examples are \cmd{USE} with \kw{ALIAS} and no alias name, \cmd{SET DELETED}, \cmd{SET HEADING}, \cmd{SET CONSOLE} or \cmd{SET STATUS} without \kw{ON} or \kw{OFF}, and \cmd{REPORT FORM} with \kw{TO FILE} and no file name.

Retype the command with the missing word supplied, and shorten any file name or condition that is too long.
\end{XbaseMessageEntry}

% 105 -- not raised: no source line sets ordinal 105 for the condition the message names (checked os/samir)
\begin{XbaseMessageEntry}[no command of this version raises this message]{105}{Table is full.}
A table of the system is full.
\end{XbaseMessageEntry}

% 53 -- not raised: no source line sets ordinal 53 for the condition the message names (checked os/samir)
\begin{XbaseMessageEntry}[no command of this version raises this message]{53}{There are no files of the type requested in this drive or catalog.}
No file of the type the command requires is on the drive or in the catalog.
\end{XbaseMessageEntry}

% 6 -- not raised: no source line sets ordinal 6 for the condition the message names (checked os/samir)
\begin{XbaseMessageEntry}[no command of this version raises this message]{6}{Too many files are open.}
More files were open at one time than the system allows.
\end{XbaseMessageEntry}

% 28 -- os/samir/samir_main.c:489 (USE: NDX_PER_AREA = 7)
\begin{XbaseMessageEntry}{28}{Too many indices.}
The \kw{INDEX} clause of \cmd{USE} named more than seven index files for one database.

Name no more than seven index files.
\end{XbaseMessageEntry}

% 74 -- not raised: no source line sets ordinal 74 for the condition the message names (checked os/samir)
\begin{XbaseMessageEntry}[no command of this version raises this message]{74}{^--- Truncated.}
A value was shortened to fit.
\end{XbaseMessageEntry}

% 92 -- os/samir/cmd/mutate.c:1343 d_run (RUN / !: the PAL could not start the command processor; os/samir/pal/pal_milton.c EXECs COMSPEC, else A:\COMMAND.COM). Gate: test-samir-dml (run)
\begin{XbaseMessageEntry}{92}{Unable to load COMMAND.COM.}
\cmd{RUN} or \cmd{!} could not load the InitechDOS command processor, \file{COMMAND.COM}: it was not found where the COMSPEC setting or drive A names it, or there was not enough memory for it.

Put the diskette with \file{COMMAND.COM} in drive A, or set COMSPEC to the drive and directory that hold it, before you start InitechBase.
\end{XbaseMessageEntry}

% 129 -- not raised: no source line sets ordinal 129 for the condition the message names (checked os/samir)
\begin{XbaseMessageEntry}[no command of this version raises this message]{129}{Unable to LOCK.}
A lock could not be taken.
\end{XbaseMessageEntry}

% 128 -- not raised: no source line sets ordinal 128 for the condition the message names (checked os/samir)
\begin{XbaseMessageEntry}[no command of this version raises this message]{128}{Unable to SKIP.}
A \cmd{SKIP} could not be carried out.
\end{XbaseMessageEntry}

% 2 -- not raised: no source line sets ordinal 2 for the condition the message names (checked os/samir)
\begin{XbaseMessageEntry}[no command of this version raises this message]{2}{Unassigned file no.}
A file number was used before a file had been assigned to it.
\end{XbaseMessageEntry}

% 133 -- not raised: no source line sets ordinal 133 for the condition the message names (checked os/samir)
\begin{XbaseMessageEntry}[no command of this version raises this message]{133}{Unauthorized access level.}
The access level is not authorized for the operation.
\end{XbaseMessageEntry}

% 144 -- not raised: no source line sets ordinal 144 for the condition the message names (checked os/samir)
\begin{XbaseMessageEntry}[no command of this version raises this message]{144}{Unauthorized duplicate.}
A duplicate key was refused.
\end{XbaseMessageEntry}

% 132 -- not raised: no source line sets ordinal 132 for the condition the message names (checked os/samir)
\begin{XbaseMessageEntry}[no command of this version raises this message]{132}{Unauthorized login.}
A login was not authorized.
\end{XbaseMessageEntry}

% 8 -- os/samir/cmd/workarea.c:1402 (g_parse_msg, the unbalanced-parenthesis parse code via xb_interp_eval); os/samir/cmd/set.c:648 (SET ALTERNATE TO: the file cannot be created). Observed in the replay: ? (1 + 2
\begin{XbaseMessageEntry}{8}{Unbalanced parenthesis.}
An expression has an opening parenthesis without a closing one.

Count the parentheses of the expression; every one that is opened must be closed.
\end{XbaseMessageEntry}

% 65 -- not raised: no source line sets ordinal 65 for the condition the message names (checked os/samir)
\begin{XbaseMessageEntry}[no command of this version raises this message]{65}{Unknown command code:}
A command code that the system does not know was met.
\end{XbaseMessageEntry}

% 104 -- not raised: no source line sets ordinal 104 for the condition the message names (checked os/samir)
\begin{XbaseMessageEntry}[no command of this version raises this message]{104}{Unknown function key.}
A function key name is not known.
\end{XbaseMessageEntry}

% 84 -- not raised: no source line sets ordinal 84 for the condition the message names (checked os/samir)
\begin{XbaseMessageEntry}[no command of this version raises this message]{84}{Unlink of old name incomplete, errno =}
Renaming a file left the file under its old name.
\end{XbaseMessageEntry}

% 16 -- os/samir/cmd/flow.c:1033 (no module accepts the verb); os/samir/samir_main.c:807,829,842 (default); os/samir/cmd/proc.c:638 (DO of an unknown procedure)
\begin{XbaseMessageEntry}{16}{*** Unrecognized command verb.}
The first word of the line is not the name of a command that InitechBase provides. \cmd{DO} reports the message when the name given is neither a procedure in the program nor a program file on the diskette.

Retype the command. Check the spelling of the verb and of the name given to \cmd{DO}.
\end{XbaseMessageEntry}

% 36 -- os/samir/samir_main.c:499 (USE clause word); os/samir/cmd/report.c:433,437,448,462
\begin{XbaseMessageEntry}{36}{Unrecognized phrase/keyword in command.}
A command contained a word that is not one of its clauses. \cmd{USE} reports the message for a word after the file name that is not \kw{ALIAS} or \kw{INDEX}; \cmd{REPORT FORM} reports it for an unknown clause, for a word that is not \kw{FILE} after \kw{TO}, and when the form name is replaced by a question mark.

Retype the command using only the clauses shown for it in this volume.
\end{XbaseMessageEntry}

% 136 -- not raised: no source line sets ordinal 136 for the condition the message names (checked os/samir)
\begin{XbaseMessageEntry}[no command of this version raises this message]{136}{Unsupported path given.}
A path given in a file name is not supported.
\end{XbaseMessageEntry}

% 35 -- os/samir/cmd/workarea.c:1401 (g_lex_msg, the lexer's unterminated-string code via xb_interp_eval). Observed in the replay: ? 'abc
\begin{XbaseMessageEntry}{35}{Unterminated string.}
A character string in an expression has no closing delimiter.

End the string with the same delimiter it begins with: an apostrophe, a quotation mark, or the bracket \texttt{]} for \texttt{[}.
\end{XbaseMessageEntry}

% 95 -- os/samir/cmd/flow.c:1466 fl_catalog (-INTERP_ERR_STRUCT typed at the dot prompt). Observed in the replay: IF .T. and ENDDO at the dot prompt
\begin{XbaseMessageEntry}{95}{Valid only in programs.}
A command that is used only in programs, such as \cmd{IF}, \cmd{ELSE}, \cmd{ENDIF}, \cmd{DO WHILE}, \cmd{ENDDO} or \cmd{LOOP}, was typed at the dot prompt.

Put the command in a program file and run the program with \cmd{DO}.
\end{XbaseMessageEntry}

% 12 -- os/samir/core/eval.c:759 (XBEE_VAR_NOT_FOUND)
\begin{XbaseMessageEntry}{12}{Variable not found.}
An expression, or a \kw{FOR} or \kw{WHILE} condition, used a name that is not a field of the database in use and not a memory variable.

Correct the spelling, open the database that holds the field with \cmd{USE}, or create the memory variable with \cmd{STORE} before it is used.
\end{XbaseMessageEntry}

% 70 -- not raised: no source line sets ordinal 70 for the condition the message names (checked os/samir)
\begin{XbaseMessageEntry}[no command of this version raises this message]{70}{** WARNING ** Data will probably be lost.  Confirm? (Y/N)}
The system asks the user to confirm an operation that may destroy data.
\end{XbaseMessageEntry}

% 94 -- os/samir/cmd/proc.c:740,835
\begin{XbaseMessageEntry}{94}{Wrong number of parameters.}
The \kw{WITH} list of \cmd{DO} held more than sixteen expressions.

Pass no more than sixteen values.
\end{XbaseMessageEntry}
